% FILE: chapters/ch06_s2_computer_vision.tex
\section{Computer Vision}\label{sec:computer_vision}

In Chapter~\ref{ch:deep_learning}, we studied the mathematical mechanics of Convolutional Neural Networks (CNNs) and Vision Transformers: sliding kernels, pooling, spatial feature hierarchies, and attention maps. In those early foundational examples, our task was primarily \textbf{image classification}: assigning a single categorical label to an entire image (for example, determining whether an inspection photo contains a ``defective'' or ``healthy'' machine component).

However, in real-world Internet of Things (IoT) systems, autonomous mobile robots (AMRs), automated optical inspection (AOI), and edge video analytics, a single global label is rarely sufficient:
\begin{itemize}
    \item A robotic fruit-harvesting arm cannot pick an apple simply by knowing that an apple exists somewhere in its camera's field of view. It needs to know the exact spatial coordinates $[x, y, z]$ of the fruit's center, its width, and its boundary to guide its pneumatic gripper without bruising the branch.
    \item An automated industrial sorting line inspecting printed circuit boards (PCBs) must identify, count, and pinpoint every missing resistor, bridge solder joint, and misaligned capacitor on a board simultaneously.
    \item An autonomous warehouse drone navigating between aisles must distinguish between the floor (drivable space), inventory pallets (obstacles), and human workers (dynamic hazards) on a continuous, pixel-by-pixel basis.
\end{itemize}

These advanced visual requirements give rise to the three core disciplines of modern \textbf{Computer Vision}: \textbf{Image Classification}, \textbf{Object Detection}, and \textbf{Semantic Segmentation}. In this section, we study how these tasks are formulated, explore real-time single-stage detectors (focusing on the architecture and mechanics of YOLO), examine pixel-level segmentation with U-Net, and learn how to optimize visual perception pipelines for microcontrollers and edge neural accelerators.

% ==============================================================================
% 1. THE COMPUTER VISION SPECTRUM
% ==============================================================================
\subsection{The Vision Spectrum: Classification, Detection, and Segmentation}\label{subsec:vision_spectrum}

Computer vision tasks exist on a spectrum of increasing spatial granularity:

\begin{figure}[htbp]
\centering
\begin{tikzpicture}[
    scale=0.88,
    box/.style={rectangle, draw=gray!70, fill=gray!5, rounded corners=2mm, thick, minimum width=3.4cm, minimum height=2.2cm, align=center},
    bbox/.style={rectangle, draw=warningred, thick, dashed, minimum width=1.6cm, minimum height=1.1cm},
    seg/.style={ellipse, draw=analogygreen!80, fill=analogygreen!20, thick, minimum width=1.4cm, minimum height=0.9cm}
]
    % 1. Classification
    \node[box] (c1) at (0, 0) {};
    \node[draw=keyideablue, fill=keyideablue!20, circle, minimum size=0.9cm] at (0, 0.2) {};
    \node[font=\footnotesize\bfseries\sffamily, text=keyideablue] at (0, -0.6) {Class: Motor};
    \node[above=0.15cm of c1, font=\small\bfseries\sffamily] {1. Classification};

    % 2. Object Detection
    \node[box] (c2) at (4.5, 0) {};
    \node[draw=keyideablue, fill=keyideablue!20, circle, minimum size=0.7cm] at (4.0, 0.2) {};
    \node[draw=iotpurple, fill=iotpurple!20, rectangle, minimum size=0.6cm] at (5.0, -0.1) {};
    \node[bbox, minimum width=1.1cm, minimum height=1.1cm] at (4.0, 0.2) {};
    \node[bbox, draw=iotpurple, minimum width=0.9cm, minimum height=0.9cm] at (5.0, -0.1) {};
    \node[font=\tiny\sffamily, text=warningred] at (4.0, 0.9) {Motor: 0.95};
    \node[font=\tiny\sffamily, text=iotpurple] at (5.0, -0.7) {Valve: 0.88};
    \node[above=0.15cm of c2, font=\small\bfseries\sffamily] {2. Object Detection};

    % 3. Semantic Segmentation
    \node[box] (c3) at (9.0, 0) {};
    \node[seg] at (8.7, 0.2) {};
    \node[ellipse, draw=iotpurple!80, fill=iotpurple!20, thick, minimum width=1.1cm, minimum height=0.8cm] at (9.4, -0.2) {};
    \node[font=\footnotesize\bfseries\sffamily, text=analogygreen!80!black] at (9.0, -0.8) {Pixel-Level Masks};
    \node[above=0.15cm of c3, font=\small\bfseries\sffamily] {3. Segmentation};
\end{tikzpicture}
\caption{The three primary tiers of computer vision: single-label classification (left), multi-object bounding box detection (center), and fine pixel-level mask segmentation (right).}\label{fig:cv_spectrum}
\end{figure}

\begin{keyidea}[Granularity of Visual Predictions]
\begin{enumerate}
    \item \textbf{Image Classification}: Predicts \emph{what} is in the image (one categorical class label for the whole image).
    \item \textbf{Object Detection}: Predicts \emph{what} and \emph{where} (multiple bounding boxes, each with a class label and confidence score).
    \item \textbf{Semantic Segmentation}: Predicts the category of \emph{every individual pixel} (dense pixel classification, creating boundary masks).
\end{enumerate}
\end{keyidea}

\paragraph{Classification on Edge Devices}
Standard deep CNNs (like ResNet-50) have over 25 million parameters and require billions of floating-point operations (FLOPs) per image, making them unviable for sub-watt microcontrollers. 

For edge classification, engineers rely on \textbf{MobileNet} architectures (Howard et al., Google). MobileNet replaces standard $3 \times 3$ convolutions with \textbf{Depthwise Separable Convolutions}:
\begin{enumerate}
    \item \textbf{Depthwise Convolution}: A single $3 \times 3$ filter is applied to each input channel separately (capturing spatial features without mixing channels).
    \item \textbf{Pointwise Convolution}: A $1 \times 1$ convolution is applied across all channels (mixing channel information).
\end{enumerate}
This factorization slashes computational cost and parameters by \textbf{$8\times$ to $9\times$} with only a minor $1\%$ drop in top-1 accuracy!

% ==============================================================================
% 2. OBJECT DETECTION: FROM TWO-STAGE TO ONE-STAGE
% ==============================================================================
\subsection{Object Detection: Two-Stage vs One-Stage Detectors}\label{subsec:detection_paradigms}

\paragraph{Why Object Detection is Hard}
In classification, the output is a fixed-size vector of class probabilities (e.g., 10 classes). In object detection:
\begin{itemize}
    \item An image may contain \textbf{zero, one, or thirty} objects. The output size is dynamic.
    \item Objects appear at \textbf{unpredictable aspect ratios and scales} (a bolt can be close to the lens or far away).
    \item The model must simultaneously solve two distinct mathematical tasks: \textbf{classification} (what is the object?) and \textbf{regression} (where are the box corners?).
\end{itemize}

\paragraph{Two-Stage Detectors (R-CNN Family)}
Early deep object detectors (Faster R-CNN, Mask R-CNN) decoupled the problem into two sequential steps:
\begin{enumerate}
    \item \textbf{Stage 1: Region Proposal}: A Region Proposal Network (RPN) scans feature maps and identifies roughly 300 to 2,000 candidate bounding boxes that likely contain objects.
    \item \textbf{Stage 2: Classification and Refinement}: Each proposed crop is cropped, resized, and passed through a classification and bounding-box regression network.
\end{enumerate}
While two-stage detectors achieve high accuracy, running hundreds of sub-crops through neural networks is computationally heavy: they operate at 5 to 15 frames per second (FPS) on expensive desktop GPUs and cannot run on edge microcontrollers.

\paragraph{One-Stage Detectors (YOLO and SSD)}
To achieve true real-time performance on constrained hardware, Joseph Redmon, Santosh Divvala, Ross Girshick, and Ali Farhadi introduced \textbf{YOLO (You Only Look Once)} in 2016. 

Instead of generating candidate proposals and evaluating each one separately, one-stage detectors frame object detection as a \textbf{single regression problem}. A single neural network takes the full input image, passes it through convolutional layers once, and predicts bounding boxes and class probabilities across all spatial locations simultaneously.

% ==============================================================================
% 3. THE YOLO ARCHITECTURE IN DEPTH
% ==============================================================================
\subsection{The YOLO (You Only Look Once) Architecture}\label{subsec:yolo}

\begin{analogy}[The Panoramic Factory Inspector]
Imagine two quality-control inspectors inspecting an electronics assembly line:
\begin{itemize}
    \item \textbf{The Two-Stage Inspector}: Uses a magnifying glass to search for suspected areas. They mark 50 candidate spots with post-it notes. Then they walk up to each post-it note, examine the component, and write down an assessment. It takes 2 minutes per board.
    \item \textbf{The YOLO Inspector}: Stands back and views the entire board in a single glance. They mentally overlay an invisible $7 \times 7$ grid across the board. In one fast visual sweep, they immediately mark every defective component, its location, and its category. The inspection takes 30 milliseconds.
\end{itemize}
\end{analogy}

\paragraph{The Grid Cell Formulation}
YOLO divides the input image into an $S \times S$ grid (e.g., $S = 7, 13$, or $19$).
\begin{itemize}
    \item If the physical center of an object falls inside a specific grid cell, \textbf{that specific grid cell is responsible for detecting that object}.
    \item Each grid cell predicts $B$ bounding boxes (typically $B = 2$ or $3$ anchor shapes).
    \item For each bounding box, the network outputs five scalar predictions:
    \begin{equation}
    \mathbf{b} = [t_x, t_y, t_w, t_h, c]\label{eq:yolo_box_tuple}
    \end{equation}
    where:
    \begin{itemize}
        \item $(t_x, t_y)$ are the coordinates of the center of the box, relative to the bounds of the grid cell (normalized to $[0, 1]$ via the sigmoid function $\sigma$).
        \item $(t_w, t_h)$ are the width and height of the box, relative to the full image dimensions (or exponential offsets relative to pre-defined anchor boxes).
        \item $c \in [0, 1]$ is the \textbf{Objectness Confidence Score}. It reflects how confident the model is that an object exists in that box, multiplied by the predicted overlap with the ground truth:
        \begin{equation}
        c = P(\text{Object}) \times \operatorname{IoU}_{\text{pred}}^{\text{truth}}\label{eq:objectness}
        \end{equation}
    \end{itemize}
    \item In addition, each cell predicts $C$ conditional class probabilities $P(\text{Class}_k \mid \text{Object})$.
\end{itemize}

The final output tensor of the network has shape:
\begin{equation}
\mathbf{Y} \in \mathbb{R}^{S \times S \times (B \times 5 + C)}\label{eq:yolo_tensor_shape}
\end{equation}
For a $7 \times 7$ grid, $B = 2$ boxes, and $C = 20$ classes (PASCAL VOC), the output tensor is $7 \times 7 \times (2 \times 5 + 20) = 7 \times 7 \times 30 = 1{,}470$ numbers. The entire image is evaluated in one single forward pass!

\begin{figure}[htbp]
\centering
\begin{tikzpicture}[
    scale=0.85,
    grid_cell/.style={rectangle, draw=gray!50, minimum size=0.9cm},
    active_cell/.style={rectangle, draw=keyideablue, fill=keyideablue!20, thick, minimum size=0.9cm},
    bbox_pred/.style={rectangle, draw=warningred, thick, rounded corners=1mm}
]
    % 7x7 Grid
    \node[font=\bfseries\sffamily] at (3.15, 3.8) {Input Image with $S \times S$ Grid ($7 \times 7$)};
    \foreach \r in {0,1,2,3,4,5,6} {
        \foreach \c in {0,1,2,3,4,5,6} {
            \node[grid_cell] (g\r\c) at (\c*0.9, 3.15 - \r*0.9) {};
        }
    }
    
    % Highlight active cell containing object center
    \node[active_cell] at (g23) {};
    \fill[warningred] (g23.center) circle (2pt) node[above right, font=\tiny\bfseries, text=warningred] {Center $(x, y)$};

    % Draw object outline (e.g. industrial motor)
    \draw[thick, draw=analogygreen!80, fill=analogygreen!10] (1.6, 1.8) rectangle (4.2, 0.4);
    \node[font=\scriptsize\bfseries\sffamily, text=analogygreen!80!black] at (2.9, 1.5) {Ground Truth Object};

    % Draw predicted bounding box
    \draw[thick, draw=warningred, dashed] (1.4, 2.0) rectangle (4.4, 0.2);
    \node[font=\tiny\bfseries\sffamily, text=warningred] at (2.9, 0.0) {Predicted Box: $[x, y, w, h, c]$};

    % Arrow pointing to output vector
    \draw[-{Stealth[length=2.5mm]}, thick, keyideablue] (6.5, 0.45) -- (8.2, 0.45);

    % Output vector representation
    \node[font=\bfseries\sffamily] at (9.8, 2.2) {Cell Output Vector};
    \draw[thick, draw=gray!70, fill=gray!5] (8.8, -1.8) rectangle (10.8, 1.8);
    \node[font=\scriptsize\sffamily] at (9.8, 1.4) {$t_x$ (Center X)};
    \node[font=\scriptsize\sffamily] at (9.8, 0.9) {$t_y$ (Center Y)};
    \node[font=\scriptsize\sffamily] at (9.8, 0.4) {$t_w$ (Width)};
    \node[font=\scriptsize\sffamily] at (9.8, -0.1) {$t_h$ (Height)};
    \node[font=\scriptsize\sffamily, text=warningred] at (9.8, -0.6) {$c$ (Confidence)};
    \node[font=\scriptsize\sffamily] at (9.8, -1.2) {$P(\text{Class}_1)$};
    \node[font=\scriptsize\sffamily] at (9.8, -1.5) {$\dots$};
\end{tikzpicture}
\caption{The YOLO Grid Formulation. The grid cell containing the object center predicts bounding box offsets, confidence, and class probabilities.}\label{fig:yolo_grid}
\end{figure}

\paragraph{Intersection over Union (IoU)}
To measure how accurately a predicted bounding box $\mathbf{B}_{\text{pred}}$ matches a ground-truth box $\mathbf{B}_{\text{truth}}$, we use \textbf{Intersection over Union (IoU)} (also known as the Jaccard Index):
\begin{equation}
\operatorname{IoU}(\mathbf{B}_{\text{pred}}, \mathbf{B}_{\text{truth}}) = \frac{\operatorname{Area}(\mathbf{B}_{\text{pred}} \cap \mathbf{B}_{\text{truth}})}{\operatorname{Area}(\mathbf{B}_{\text{pred}} \cup \mathbf{B}_{\text{truth}})}\label{eq:iou_formula}
\end{equation}
where $\operatorname{Area}(\mathbf{B}_{\text{pred}} \cup \mathbf{B}_{\text{truth}}) = \operatorname{Area}(\mathbf{B}_{\text{pred}}) + \operatorname{Area}(\mathbf{B}_{\text{truth}}) - \operatorname{Area}(\mathbf{B}_{\text{pred}} \cap \mathbf{B}_{\text{truth}})$.
\begin{itemize}
    \item $\operatorname{IoU} = 1.0$: Perfect, exact overlap.
    \item $\operatorname{IoU} \ge 0.50$: Standard historical benchmark for a ``correct'' detection.
    \item $\operatorname{IoU} \ge 0.75$: Strict benchmark for high-precision robotic manipulation.
\end{itemize}

\begin{example}[Hand-Calculating Intersection over Union (IoU)]
Let two bounding boxes be defined by their corner coordinates $[x_{\min}, y_{\min}, x_{\max}, y_{\max}]$ on a $100 \times 100$ image:
\begin{align*}
\mathbf{B}_{\text{truth}} &= [20, 20, 60, 60] \quad (\text{Width} = 40, \text{Height} = 40) \\
\mathbf{B}_{\text{pred}} &= [30, 20, 70, 60] \quad (\text{Width} = 40, \text{Height} = 40)
\end{align*}

\textbf{Step 1: Compute Individual Box Areas}
\begin{align*}
\operatorname{Area}(\mathbf{B}_{\text{truth}}) &= (60 - 20) \times (60 - 20) = 40 \times 40 = \mathbf{1{,}600} \\
\operatorname{Area}(\mathbf{B}_{\text{pred}}) &= (70 - 30) \times (60 - 20) = 40 \times 40 = \mathbf{1{,}600}
\end{align*}

\textbf{Step 2: Find the Intersection Rectangle}
The intersection boundaries are defined by:
\begin{align*}
x_{\min}^{\text{inter}} &= \max(x_{\min}^{\text{truth}}, x_{\min}^{\text{pred}}) = \max(20, 30) = 30 \\
y_{\min}^{\text{inter}} &= \max(y_{\min}^{\text{truth}}, y_{\min}^{\text{pred}}) = \max(20, 20) = 20 \\
x_{\max}^{\text{inter}} &= \min(x_{\max}^{\text{truth}}, x_{\max}^{\text{pred}}) = \min(60, 70) = 60 \\
y_{\max}^{\text{inter}} &= \min(y_{\max}^{\text{truth}}, y_{\max}^{\text{pred}}) = \min(60, 60) = 60
\end{align*}
Compute intersection dimensions and area:
\begin{align*}
\text{Width}_{\text{inter}} &= \max(0, 60 - 30) = 30 \\
\text{Height}_{\text{inter}} &= \max(0, 60 - 20) = 40 \\
\operatorname{Area}(\text{Intersection}) &= 30 \times 40 = \mathbf{1{,}200}
\end{align*}

\textbf{Step 3: Compute Union Area}
\begin{align*}
\operatorname{Area}(\text{Union}) &= \operatorname{Area}(\mathbf{B}_{\text{truth}}) + \operatorname{Area}(\mathbf{B}_{\text{pred}}) - \operatorname{Area}(\text{Intersection}) \\
                                  &= 1{,}600 + 1{,}600 - 1{,}200 = 3{,}200 - 1{,}200 = \mathbf{2{,}000}
\end{align*}

\textbf{Step 4: Compute IoU}
\begin{equation*}
\operatorname{IoU} = \frac{1{,}200}{2{,}000} = \frac{12}{20} = \mathbf{0.60} \quad (60\%)
\end{equation*}
Since $\operatorname{IoU} = 0.60 \ge 0.50$, this detection satisfies the standard PASCAL VOC criterion for a successful true positive.
\end{example}

\paragraph{Post-Processing: Non-Maximum Suppression (NMS)}
Because multiple adjacent grid cells often detect the same physical object with slight offsets, a raw YOLO forward pass produces multiple overlapping bounding boxes around each object. To prune duplicates, we run \textbf{Non-Maximum Suppression (NMS)}:
\begin{enumerate}
    \item Discard all boxes with confidence score below a detection threshold (e.g., $c < 0.25$).
    \item Select the remaining box with the highest confidence score and save it to the final detection list.
    \item Calculate the IoU between this selected box and all remaining candidate boxes.
    \item Discard any candidate box that has an $\operatorname{IoU} > \text{threshold}_{\text{NMS}}$ (typically $0.45$).
    \item Repeat steps 2--4 until no candidate boxes remain.
\end{enumerate}

% ==============================================================================
% 4. SEGMENTATION BASICS: U-NET AND PIXEL MASKS
% ==============================================================================
\subsection{Segmentation Basics: U-Net and Pixel-Level Classification}\label{subsec:segmentation}

While object detection draws rectangular bounding boxes, many edge applications require knowing the \textbf{exact physical boundary} of an object. For instance, a drone inspecting a concrete wind-turbine tower must measure the exact square millimeters of surface crack erosion, not just draw a loose box around it.

\paragraph{Semantic vs Instance Segmentation}
\begin{itemize}
    \item \textbf{Semantic Segmentation}: Assigns every pixel a category label (e.g., ``crack'', ``rust'', ``background''). It does not differentiate between separate instances of the same object.
    \item \textbf{Instance Segmentation}: Distinguishes between separate instances (e.g., labeling Crack \#1 and Crack \#2 as distinct entities).
\end{itemize}

\paragraph{The U-Net Architecture}
Introduced in 2015 by Olaf Ronneberger, Philipp Fischer, and Thomas Brox for biomedical image segmentation, the \textbf{U-Net} is the premier architecture for dense pixel-level prediction.

\begin{figure}[htbp]
\centering
\begin{tikzpicture}[
    scale=0.82,
    enc_block/.style={rectangle, draw=keyideablue, fill=keyideablue!20, thick, minimum width=1.4cm, minimum height=0.7cm, font=\scriptsize\sffamily\bfseries},
    dec_block/.style={rectangle, draw=analogygreen, fill=analogygreen!20, thick, minimum width=1.4cm, minimum height=0.7cm, font=\scriptsize\sffamily\bfseries},
    pool/.style={-{Stealth[length=2.5mm]}, draw=warningred, thick},
    up/.style={-{Stealth[length=2.5mm]}, draw=analogygreen!80!black, thick},
    skip/.style={-{Stealth[length=2mm]}, draw=keyideablue!80, dashed, thick}
]
    % Contracting Path (Encoder - Left side)
    \node[enc_block] (e1) at (0, 3) {Conv $3 \times 3$\\($H \times W$)};
    \node[enc_block] (e2) at (1.8, 1.8) {Conv $3 \times 3$\\($\frac{H}{2} \times \frac{W}{2}$)};
    \node[enc_block] (e3) at (3.6, 0.6) {Conv $3 \times 3$\\($\frac{H}{4} \times \frac{W}{4}$)};
    \node[enc_block, fill=iotpurple!20, draw=iotpurple] (b) at (5.4, -0.6) {Bottleneck\\($\frac{H}{8} \times \frac{W}{8}$)};

    % Expanding Path (Decoder - Right side)
    \node[dec_block] (d3) at (7.2, 0.6) {UpConv\\($\frac{H}{4} \times \frac{W}{4}$)};
    \node[dec_block] (d2) at (9.0, 1.8) {UpConv\\($\frac{H}{2} \times \frac{W}{2}$)};
    \node[dec_block] (d1) at (10.8, 3) {UpConv\\($H \times W$)};

    % Downsampling connections (MaxPool)
    \draw[pool] (e1.south) |- (e2.west) node[near start, left, font=\tiny\text=warningred] {MaxPool};
    \draw[pool] (e2.south) |- (e3.west) node[near start, left, font=\tiny\text=warningred] {MaxPool};
    \draw[pool] (e3.south) |- (b.west) node[near start, left, font=\tiny\text=warningred] {MaxPool};

    % Upsampling connections
    \draw[up] (b.east) -| (d3.south) node[near start, right, font=\tiny\text=analogygreen] {UpSample};
    \draw[up] (d3.east) -| (d2.south) node[near start, right, font=\tiny\text=analogygreen] {UpSample};
    \draw[up] (d2.east) -| (d1.south) node[near start, right, font=\tiny\text=analogygreen] {UpSample};

    % Skip Connections
    \draw[skip] (e1) -- node[above, font=\tiny\bfseries] {Skip Connection (Copy \& Concat)} (d1);
    \draw[skip] (e2) -- (d2);
    \draw[skip] (e3) -- (d3);

    % Input and Output
    \node[left=0.6cm of e1, font=\footnotesize\bfseries\sffamily, align=center] {Input\\Image};
    \node[right=0.6cm of d1, font=\footnotesize\bfseries\sffamily, align=center] {Pixel Mask\\Prediction};
\end{tikzpicture}
\caption{The U-Net Architecture. The encoder extracts semantic context; the decoder restores spatial resolution; skip connections pass high-resolution spatial details directly across.}\label{fig:unet}
\end{figure}

\begin{keyidea}[The Magic of Skip Connections in U-Net]
In a standard encoder-decoder network, repeated pooling operations downsample the image, destroying fine spatial details (such as thin hairline cracks or exact object boundaries). 

U-Net solves this by adding \textbf{long skip connections} that copy high-resolution feature maps directly from each encoder level across to the corresponding decoder level. The decoder combines high-level semantic context (``this is a crack'') with sharp low-level spatial geometry (``the boundary sits exactly at pixel row 42'').
\end{keyidea}

\paragraph{Loss Functions for Segmentation: The Class Imbalance Problem}
In surface defect detection, a crack may occupy only $100$ pixels on a $1{,}000 \times 1{,}000$ image ($0.01\%$ of the total pixels). If we train with standard Binary Cross-Entropy (BCE), a model that predicts ``healthy background'' for every single pixel achieves \textbf{$99.99\%$ accuracy} while being completely useless!

To solve this extreme imbalance, segmentation models use the \textbf{Dice Loss}:
\begin{equation}
\mathcal{L}_{\text{Dice}} = 1 - \frac{2 \sum_{i=1}^N p_i y_i + \epsilon}{\sum_{i=1}^N p_i + \sum_{i=1}^N y_i + \epsilon}\label{eq:dice_loss}
\end{equation}
where $p_i \in [0, 1]$ is the predicted probability for pixel $i$, $y_i \in \{0, 1\}$ is the ground-truth mask, and $\epsilon$ prevents division by zero. Dice Loss directly optimizes the overlap (F1-score) of the target mask, ignoring the overwhelming sea of background pixels.

% ==============================================================================
% 5. CODE WALKTHROUGH: IOU AND NMS IN PYTORCH
% ==============================================================================
\subsection{PyTorch Implementation: IoU and Non-Maximum Suppression}\label{subsec:pytorch_nms}

Let us build a complete, self-contained PyTorch script implementing the core algorithms that power every real-time object detector: \textbf{Vectorized Intersection over Union (IoU)} and \textbf{Non-Maximum Suppression (NMS)}.

\paragraph{The PyTorch Implementation}
\begin{lstlisting}[language=Python]
import torch

def compute_iou(boxes1, boxes2):
    """
    Computes pairwise Intersection over Union (IoU) between two sets of boxes.
    Boxes format: [x1, y1, x2, y2]
    """
    # 1. Determine intersection coordinates
    x1 = torch.max(boxes1[:, 0].unsqueeze(1), boxes2[:, 0].unsqueeze(0))
    y1 = torch.max(boxes1[:, 1].unsqueeze(1), boxes2[:, 1].unsqueeze(0))
    x2 = torch.min(boxes1[:, 2].unsqueeze(1), boxes2[:, 2].unsqueeze(0))
    y2 = torch.min(boxes1[:, 3].unsqueeze(1), boxes2[:, 3].unsqueeze(0))
    
    # 2. Calculate intersection area
    intersection = torch.clamp(x2 - x1, min=0) * torch.clamp(y2 - y1, min=0)
    
    # 3. Calculate individual box areas
    area1 = (boxes1[:, 2] - boxes1[:, 0]) * (boxes1[:, 3] - boxes1[:, 1])
    area2 = (boxes2[:, 2] - boxes2[:, 0]) * (boxes2[:, 3] - boxes2[:, 1])
    
    # 4. Compute union and IoU
    union = area1.unsqueeze(1) + area2.unsqueeze(0) - intersection
    iou = intersection / torch.clamp(union, min=1e-6)
    return iou

def non_max_suppression(boxes, scores, iou_threshold=0.5):
    """
    Applies Non-Maximum Suppression (NMS) to eliminate overlapping redundant boxes.
    """
    # Sort boxes in descending order of confidence score
    sorted_indices = torch.argsort(scores, descending=True)
    keep_boxes = []
    
    while sorted_indices.numel() > 0:
        # Pick the box with the highest remaining score
        current = sorted_indices[0].item()
        keep_boxes.append(current)
        
        if sorted_indices.numel() == 1:
            break
            
        # Compute IoU between current highest box and all other remaining boxes
        current_box = boxes[current].unsqueeze(0)
        remaining_boxes = boxes[sorted_indices[1:]]
        ious = compute_iou(current_box, remaining_boxes).squeeze(0)
        
        # Keep only candidate boxes whose IoU is BELOW the threshold
        below_thresh_mask = ious < iou_threshold
        # Offset indices by 1 because we sliced sorted_indices[1:]
        sorted_indices = sorted_indices[1:][below_thresh_mask]
        
    return keep_boxes

# Verification Test: Simulated detections on a conveyor belt
# Three bounding boxes: Box 0 and Box 1 overlap heavily (same part); Box 2 is separate
boxes = torch.tensor([
    [20.0, 20.0, 60.0, 60.0],  # Box 0: Confidence 0.92
    [22.0, 21.0, 61.0, 59.0],  # Box 1: Confidence 0.85 (duplicate of Box 0)
    [75.0, 70.0, 95.0, 90.0]   # Box 2: Confidence 0.88 (distinct component)
])
scores = torch.tensor([0.92, 0.85, 0.88])

# Pairwise IoU between Box 0 and Box 1
pair_iou = compute_iou(boxes[0:1], boxes[1:2]).item()
print(f"Pairwise IoU between Box 0 and Box 1: {pair_iou:.4f}")

# Execute NMS
kept_indices = non_max_suppression(boxes, scores, iou_threshold=0.5)
print(f"Indices kept after NMS: {kept_indices}")
for idx in kept_indices:
    print(f"  Kept Box {idx}: {boxes[idx].tolist()} with Score {scores[idx]:.2f}")
\end{lstlisting}

\paragraph{Expected Output}
\begin{lstlisting}[style=outputstyle]
Pairwise IoU between Box 0 and Box 1: 0.8647
Indices kept after NMS: [0, 2]
  Kept Box 0: [20.0, 20.0, 60.0, 60.0] with Score 0.92
  Kept Box 2: [75.0, 70.0, 95.0, 90.0] with Score 0.88
\end{lstlisting}

\paragraph{Line-by-Line Code Breakdown}
\begin{itemize}
    \item \textbf{Lines 8--11}: \code{x1 = torch.max(boxes1[:, 0].unsqueeze(1), boxes2[:, 0].unsqueeze(0))}. Uses PyTorch broadcasting to compute the pairwise intersection corners for all pairs of boxes simultaneously without Python \code{for} loops.
    \item \textbf{Line 14}: \code{torch.clamp(x2 - x1, min=0)} ensures that if two boxes do not overlap, their intersection length evaluates to $0.0$ rather than a negative number.
    \item \textbf{Line 30}: \code{sorted_indices = torch.argsort(scores, descending=True)}. Orders the candidate boxes from most confident to least confident.
    \item \textbf{Lines 43--45}: \code{ious < iou_threshold}. Discards Box 1 because its overlap with Box 0 is $86.5\%$ (exceeding the $0.50$ threshold). Box 2 is preserved because its overlap is $0.0\%$.
\end{itemize}

% ==============================================================================
% 6. IOT & TINYML CONNECTION
% ==============================================================================
\subsection{ML for IoT and Edge: Embedded Vision on Microcontrollers}\label{subsec:iot_vision}

\begin{iotbox}[The Edge Vision Reality: Hardware vs Frame Rate]
Running real-time computer vision on embedded hardware involves choosing the right compute tier:
\begin{itemize}
    \item \textbf{Tier 1: Microcontrollers (ESP32-CAM, STM32H7, Raspberry Pi Pico)}:
    \begin{itemize}
        \item 256 KB to 8 MB RAM, 160 to 480 MHz CPU, no GPU.
        \item Can run \textbf{FOMO} or compact MobileNet classifiers at $5\text{--}15\text{ FPS}$ on $96 \times 96$ grayscale images.
        \item Cannot run full YOLO architectures due to SRAM memory limits.
    \end{itemize}
    \item \textbf{Tier 2: Edge Neural Accelerators (Kendryte K210, Coral Edge TPU, Arm Ethos)}:
    \begin{itemize}
        \item Dedicated hardware NPU ($0.5$ to $4\text{ TOPS}$, $<2\text{ Watts}$).
        \item Runs \textbf{Tiny-YOLO} or \textbf{MobileNet-SSD} at $30\text{--}60\text{ FPS}$ on $320 \times 240$ color video.
    \end{itemize}
\end{itemize}
\end{iotbox}

\paragraph{FOMO: Faster Objects, More Objects on Microcontrollers}
Developed by Edge Impulse for microcontrollers like the ESP32, \textbf{FOMO} bridges the gap between image classification and object detection:
\begin{itemize}
    \item \textbf{The Insight}: On a microcontroller, estimating bounding box width and height $(w, h)$ consumes excessive compute and RAM. But for counting items on a conveyor belt or detecting bee entries at a beehive, \textbf{you only need the $(x, y)$ centroid coordinates}!
    \item FOMO takes a MobileNet backbone, removes the dense classification head, and outputs a coarse grid of class probabilities.
    \item Instead of drawing boxes, FOMO marks centroids. It uses \textbf{$30\times$ less memory} than MobileNet-SSD, fitting into under 150 KB of RAM and running at 30 FPS on an ESP32!
\end{itemize}

% ==============================================================================
% 7. COMMON MISTAKES AND MISCONCEPTIONS
% ==============================================================================
\subsection{Common Mistakes and Misconceptions}\label{subsec:mistakes}

\begin{warning}[Five Pitfalls in Edge Computer Vision]
\begin{enumerate}
    \item \textbf{Confusing Box Coordinate Formats}: Mixing up the two standard bounding-box formats:
    \begin{itemize}
        \item \textbf{Corner format}: $[x_{\min}, y_{\min}, x_{\max}, y_{\max}]$ (used by torchvision and IoU calculation).
        \item \textbf{Center-Dimension format}: $[x_{\text{center}}, y_{\text{center}}, \text{width}, \text{height}]$ (used by YOLO).
    \end{itemize}
    Feeding center coordinates into an algorithm expecting corner coordinates produces nonsensical negative areas and inverted boxes.
    \item \textbf{Deploying Full YOLO to Raw Microcontrollers}: Attempting to flash full YOLOv5 or YOLOv8 onto an ESP32. Even the smallest YOLOv8n model has 3.2 million parameters ($\approx 12.8\text{ MB}$ in Float32) and requires over 20 MB of intermediate activation memory during inference, immediately crashing the microcontroller. Microcontrollers require dedicated TinyML architectures like FOMO or INT8 MobileNet.
    \item \textbf{Tuning NMS Thresholds Incorrectly}: Setting the NMS IoU threshold too low (e.g., $0.10$). If two distinct physical objects overlap slightly (like two apples resting side by side in a crate), an aggressive NMS threshold suppresses the second apple, leading to under-counting.
    \item \textbf{Evaluating Segmentation with Pixel Accuracy on Imbalanced Data}: Reporting $99\%$ pixel accuracy on a dataset where the defect occupies $1\%$ of the surface area. Always evaluate segmentation models using \textbf{Mean IoU (mIoU)} or the \textbf{Dice Coefficient}.
    \item \textbf{Overlooking Camera Motion Blur on Moving Conveyors}: Training a high-accuracy vision model on sharp, static photos, and then deploying it to an edge camera mounted over a vibrating, high-speed conveyor belt. Shutter speed and motion blur degrade features. Always train with motion-blur and exposure data augmentations.
\end{enumerate}
\end{warning}

% ==============================================================================
% 8. KEY TAKEAWAYS
% ==============================================================================
\subsection{Key Takeaways}\label{subsec:takeaways}

\begin{takeaways}[Core Principles of Computer Vision]
\begin{itemize}
    \item \textbf{The Computer Vision Spectrum}: Classification assigns a single global label; Object Detection identifies bounding boxes $[x, y, w, h, c]$; Segmentation classifies every individual pixel.
    \item \textbf{One-Stage Speed Advantage}: YOLO evaluates an image in a single forward pass by dividing it into an $S \times S$ grid, enabling real-time detection at 30+ FPS.
    \item \textbf{IoU and NMS}: Intersection over Union (IoU) measures spatial overlap accuracy; Non-Maximum Suppression (NMS) prunes redundant overlapping detections.
    \item \textbf{U-Net Skip Connections}: U-Net passes high-resolution spatial feature maps directly from encoder stages to decoder stages, preserving precise pixel boundaries.
    \item \textbf{TinyML Edge Vision}: On resource-constrained microcontrollers (ESP32), centroid tracking (FOMO) and depthwise separable convolutions (MobileNet) replace full bounding-box regression to satisfy strict SRAM limits.
\end{itemize}
\end{takeaways}

% ==============================================================================
% 9. CHECK YOURSELF
% ==============================================================================
\subsection{Check Yourself}\label{subsec:check_yourself}

\begin{checkyourself}[Practice Questions and Calculations]
\begin{enumerate}
    \item \textbf{Intersection over Union (IoU) Calculation}: Two bounding boxes are given in corner coordinates $[x_1, y_1, x_2, y_2]$:
    \begin{equation*}
    \mathbf{B}_A = [10, 10, 50, 50], \quad \mathbf{B}_B = [30, 10, 70, 50]
    \end{equation*}
    Calculate the exact Intersection over Union (IoU) between Box A and Box B.
    \item \textbf{YOLO Output Tensor Dimensioning}: An edge detector divides an image into a $13 \times 13$ grid. Each grid cell predicts $B = 3$ bounding boxes. The system must detect 5 distinct industrial component classes. Exactly how many scalar values are contained in the final output tensor of this network?
    \item \textbf{Non-Maximum Suppression Logic}: Three bounding boxes predict the same object class with confidences $c_1 = 0.90, c_2 = 0.80, c_3 = 0.70$. The pairwise IoUs are $\operatorname{IoU}(1, 2) = 0.65$ and $\operatorname{IoU}(1, 3) = 0.20$. If the NMS threshold is set to $\tau_{\text{NMS}} = 0.50$, which boxes are kept and which are suppressed?
    \item \textbf{Dice Coefficient Calculation}: In a $10 \times 10$ binary segmentation mask, the ground-truth crack occupies 8 pixels. The model predicts 10 positive pixels, of which 6 overlap with the ground truth. Calculate the Dice coefficient.
    \item \textbf{Depthwise Separable Convolution Efficiency}: A standard $3 \times 3$ convolution takes an input tensor with 64 channels and produces 64 output channels. A depthwise separable convolution performs a $3 \times 3$ depthwise convolution followed by a $1 \times 1$ pointwise convolution. Exactly how many weights does each method require (ignoring biases), and what is the percentage parameter reduction?
\end{enumerate}
\end{checkyourself}

% ==============================================================================
% 10. ANSWERS
% ==============================================================================
\subsection*{Answers}\label{subsec:answers}

\begin{enumerate}
    \item \textbf{Intersection over Union (IoU) Calculation}:
    \begin{itemize}
        \item Individual Box Areas:
        \begin{align*}
        \operatorname{Area}(\mathbf{B}_A) &= (50 - 10) \times (50 - 10) = 40 \times 40 = 1{,}600 \\
        \operatorname{Area}(\mathbf{B}_B) &= (70 - 30) \times (50 - 10) = 40 \times 40 = 1{,}600
        \end{align*}
        \item Intersection Coordinates:
        \begin{align*}
        x_1^{\text{inter}} &= \max(10, 30) = 30, \quad y_1^{\text{inter}} = \max(10, 10) = 10 \\
        x_2^{\text{inter}} &= \min(50, 70) = 50, \quad y_2^{\text{inter}} = \min(50, 50) = 50 \\
        \text{Width} &= 50 - 30 = 20, \quad \text{Height} = 50 - 10 = 40 \\
        \operatorname{Area}(\text{Intersection}) &= 20 \times 40 = \mathbf{800}
        \end{align*}
        \item Union Area:
        \begin{equation*}
        \operatorname{Area}(\text{Union}) = 1{,}600 + 1{,}600 - 800 = \mathbf{2{,}400}
        \end{equation*}
        \item IoU:
        \begin{equation*}
        \operatorname{IoU} = \frac{800}{2{,}400} = \frac{1}{3} \approx \mathbf{0.3333} \quad (33.3\%)
        \end{equation*}
    \end{itemize}

    \item \textbf{YOLO Output Tensor Dimensioning}:
    For each cell, the output contains:
    \begin{equation*}
    B \times 5 + C = (3 \times 5) + 5 = 15 + 5 = 20 \text{ values per cell}
    \end{equation*}
    Across the $13 \times 13$ spatial grid:
    \begin{equation*}
    \text{Total Tensor Size} = 13 \times 13 \times 20 = 169 \times 20 = \mathbf{3{,}380 \text{ scalar values}}
    \end{equation*}

    \item \textbf{Non-Maximum Suppression Logic}:
    \begin{enumerate}
        \item Step 1: Box 1 has the highest confidence ($0.90$). Box 1 is **kept**.
        \item Step 2: Compare remaining boxes to Box 1:
        \begin{itemize}
            \item $\operatorname{IoU}(1, 2) = 0.65 > 0.50 \implies$ Box 2 is **suppressed (discarded)**.
            \item $\operatorname{IoU}(1, 3) = 0.20 \le 0.50 \implies$ Box 3 is **retained**.
        \end{itemize}
        \item Step 3: Box 3 is the highest remaining box ($0.70$). Box 3 is **kept**.
        \item Final result: Boxes **1 and 3 are kept**; Box 2 is eliminated.
    \end{enumerate}

    \item \textbf{Dice Coefficient Calculation}:
    Using Equation~\eqref{eq:dice_loss}:
    \begin{equation*}
    \text{Dice} = \frac{2 |\mathbf{X} \cap \mathbf{Y}|}{|\mathbf{X}| + |\mathbf{Y}|} = \frac{2 \times 6}{8 + 10} = \frac{12}{18} = \frac{2}{3} \approx \mathbf{0.6667}
    \end{equation*}

    \item \textbf{Depthwise Separable Convolution Efficiency}:
    Let $C_{\text{in}} = 64, C_{\text{out}} = 64, K = 3$.
    \begin{itemize}
        \item \textbf{Standard Convolution}:
        \begin{equation*}
        \text{Weights} = K \times K \times C_{\text{in}} \times C_{\text{out}} = 3 \times 3 \times 64 \times 64 = 9 \times 4{,}096 = \mathbf{36{,}864 \text{ weights}}
        \end{equation*}
        \item \textbf{Depthwise Separable Convolution}:
        \begin{align*}
        \text{Depthwise Weights} &= 3 \times 3 \times C_{\text{in}} = 9 \times 64 = 576 \\
        \text{Pointwise Weights} &= 1 \times 1 \times C_{\text{in}} \times C_{\text{out}} = 1 \times 64 \times 64 = 4{,}096 \\
        \text{Total Weights} &= 576 + 4{,}096 = \mathbf{4{,}672 \text{ weights}}
        \end{align*}
        \item \textbf{Percentage Reduction}:
        \begin{equation*}
        \text{Reduction} = \frac{36{,}864 - 4{,}672}{36{,}864} = \frac{32{,}192}{36{,}864} \approx \mathbf{87.3\%}
        \end{equation*}
        Depthwise separable convolution reduces parameters and compute by over $87\%$.
    \end{itemize}
\end{enumerate}
